\chapterimage{figs/chapterhead.png}
\chapter{Biological and Whole-Cell Extensions}
\label{chap:biological-and-whole-cell-extensions}

This chapter maps the biological extension surface that exists in the current
tree. It is intentionally conservative: it documents implemented classes,
didactic examples, and validation boundaries, but it does not claim organism-
level predictive validity.

\noindent\textbf{Objective.} Summarize the biological, whole-cell, and related
extension areas supported by the project.

\noindent\textbf{Audience.} Developers and researchers who need a high-level map
of the biological extension surface.

\noindent\textbf{Scope.} Cover the scope, terminology, and evidence limits that
should govern these extensions.

\noindent\textbf{Status.} Temporary skeleton replaced with current implemented
surface and evidence boundaries. The scientific claim level remains
experimental/research-oriented unless a later chapter states otherwise.

\section{Scope and contract}

The repository currently treats whole-cell, biochemical, SBML, and virtual-cell
work as experimental or research-oriented. That boundary is stated in
\path{../README.md}, \path{../docs/ai_assistants/reference/SCIENTIFIC_DOMAINS.md},
and the human decision backlog. The right reading is:
\begin{itemize}
\item the code exposes real classes, examples, and tests;
\item the examples are didactic unless a source file explicitly says otherwise;
\item successful compilation or startup does not imply scientific validation;
\item any future claim beyond the current code must bring a real reference,
      dataset or model identifier, benchmark, tolerance, procedure, and
      explicit scope.
\end{itemize}

The biological surface in the current tree is split across three layers:
\begin{itemize}
\item native biochemical definitions and bridges;
\item whole-cell state and components;
\item model-specific didactic examples.
\end{itemize}

The key design rule is that biological meaning stays explicit. A biochemical
network, a flux-balance network, a compartment transport rule, and a whole-cell
state are not interchangeable abstractions.

The key visual references for this chapter are Figure~\ref{fig:biochemical-network-surface},
Figure~\ref{fig:fba-flow},
Figure~\ref{fig:compartment-transport-flow},
Figure~\ref{fig:whole-cell-architecture},
Figure~\ref{fig:yeast-mvp-structure}, and
Figure~\ref{fig:pathway-stress-example}.

\section{Biochemical surface}

The biochemical layer currently includes both native deterministic networks and
interoperability helpers.

\subsection{Native deterministic and network-oriented classes}

The current surface centers on the following classes:
\begin{itemize}
\item \texttt{BioSpecies}, \texttt{BioParameter}, and \texttt{BioReaction} for
      native species, scalar parameters, and mass-action reactions;
\item \texttt{BioNetwork} for a deterministic biochemical network runner with
      fixed-step RK4 over a mass-action ODE system;
\item \texttt{MetabolicNetwork} and \texttt{MetabolicReaction} for flux-style
      metabolic networks that keep reaction membership and objective metadata
      explicit;
\item \texttt{GeneticCircuit}, \texttt{GeneticRegulation}, and related genetic
      circuit helpers for transcriptional regulation surfaces;
\item \texttt{BioSimulatorRunner} as the command-driven structural runner for
      validate/simulate/report/import/export workflows;
\item \texttt{BioSBMLBridge} as the native import/export boundary for SBML text.
\end{itemize}

Figure~\ref{fig:biochemical-network-surface} sketches that surface as it exists
today: native biochemical definitions feed either the deterministic runner or
the SBML bridge, while the flux-style network types stay available for the
whole-cell examples that need them.

\begin{figure}[htbp]
\centering
% FIGURE-SPEC-BEGIN
% ID: fig:biochemical-network-surface
% Source of truth:
%   - ../source/plugins/data/BiochemicalSimulation/BioSpecies.h
%   - ../source/plugins/data/BiochemicalSimulation/BioParameter.h
%   - ../source/plugins/data/BiochemicalSimulation/BioReaction.h
%   - ../source/plugins/data/BiochemicalSimulation/BioNetwork.h
%   - ../source/plugins/data/BiochemicalSimulation/BioSimulatorRunner.h
%   - ../source/plugins/data/BiochemicalSimulation/BioSBMLBridge.h
% Purpose:
%   Show the current biochemical extension surface and its import/export boundary.
% Required visual content:
%   - native biochemical definitions
%   - deterministic network runner
%   - SBML bridge
%   - structural runner commands
% Accessibility description:
%   A block diagram that shows native biochemical data flowing into the deterministic runner and into the SBML bridge.
% Validation:
%   The diagram must match the current class names and the command-oriented runner surface in the codebase.
% FIGURE-SPEC-END
\figureplaceholder{figs/generated/biochemical-network-surface}
\caption{Current biochemical extension surface and its SBML boundary.}
\label{fig:biochemical-network-surface}
\end{figure}

\subsection{Flux-style metabolic models}

The metabolic path is deliberately separate from the deterministic network
runner. \texttt{MetabolicNetwork} stores the reaction membership, objective
reaction, objective sense, compartment, and enabled flag. \texttt{MetabolicReaction}
stores reactants, products, bounds, objective coefficient, reversibility, and
gene rule text.

This is the layer used by the whole-cell examples that exercise FBA-style
reasoning. The manual should treat it as a didactic constraint model rather than
as a validated biochemical simulator.

\begin{figure}[htbp]
\centering
% FIGURE-SPEC-BEGIN
% ID: fig:fba-flow
% Source of truth:
%   - ../source/plugins/data/BiochemicalSimulation/MetabolicNetwork.h
%   - ../source/plugins/data/BiochemicalSimulation/MetabolicReaction.h
%   - ../source/plugins/components/BiochemicalSimulation/MetabolicFluxBalance.h
%   - ../source/plugins/components/WholeCellModeling/MetabolicStateProjectionComponent.h
% Purpose:
%   Show the FBA-style path from reaction membership to flux balance and projection.
% Required visual content:
%   - metabolic reaction membership and bounds
%   - objective selection
%   - flux-balance solve
%   - projection into pathway activity or pools
% Accessibility description:
%   A pipeline diagram from a metabolic network through FBA to projected state updates.
% Validation:
%   The diagram must reflect the current backend behavior in the whole-cell MVPs and tests.
% FIGURE-SPEC-END
\figureplaceholder{figs/generated/fba-flow}
\caption{FBA-style metabolic flow used by the whole-cell examples.}
\label{fig:fba-flow}
\end{figure}

\subsection{Compartment transport}

Compartment-aware transport is modeled explicitly rather than hidden inside the
solver. The main helper classes are:
\begin{itemize}
\item \texttt{MetabolicStateProjectionComponent} for mapping selected fluxes into
      global metabolite pools, compartment pools, and pathway activity values;
\item \texttt{CompartmentExchangeComponent} for moving selected amounts between
      regions such as extracellular, cytosol, mitochondria, or bud;
\item \texttt{ResourceAllocationComponent} for distributing RNAP and ribosome
      resources before expression steps.
\end{itemize}

Figure~\ref{fig:compartment-transport-flow} captures the compartment-aware flow
used by the yeast-oriented whole-cell examples.

\begin{figure}[htbp]
\centering
% FIGURE-SPEC-BEGIN
% ID: fig:compartment-transport-flow
% Source of truth:
%   - ../source/plugins/components/WholeCellModeling/MetabolicStateProjectionComponent.h
%   - ../source/plugins/components/WholeCellModeling/CompartmentExchangeComponent.h
%   - ../source/plugins/components/WholeCellModeling/ResourceAllocationComponent.h
%   - ../source/applications/modelSpecific/wcm/WcmCompartmentTransportWholeCellMvp.h
%   - ../source/applications/modelSpecific/wcm/WcmYeastDidacticWholeCellMvp.h
% Purpose:
%   Show how fluxes become compartment pools and how exchange rules consume pathway activity.
% Required visual content:
%   - source compartment
%   - target compartment
%   - projected flux
%   - exchange rule
%   - downstream expression/growth stages
% Accessibility description:
%   A compartment map that shows transport, projection, and downstream consumption of metabolites.
% Validation:
%   The diagram must match the current whole-cell transport examples and their runtime order.
% FIGURE-SPEC-END
\figureplaceholder{figs/generated/compartment-transport-flow}
\caption{Compartment transport and projection in the current whole-cell examples.}
\label{fig:compartment-transport-flow}
\end{figure}

\section{Whole-cell surface}

The whole-cell layer centers on \texttt{WholeCellState}. That class is the
shared container used by the whole-cell components. Its current responsibilities
include:
\begin{itemize}
\item integer molecule counts for SSA-driven expression and degradation;
\item metabolite pools, including compartment-scoped pools;
\item pathway activity values used by growth, checkpoint, and stress routing;
\item resource budgets for expression allocation;
\item cell volume, mass, time, step count, lifecycle phase, generation count,
      viability, and last-division time;
\item JSON path storage for the CovertLab/WholeCell fixed-constants and
      parameters inputs used by the didactic M. genitalium-oriented path.
\end{itemize}

The main whole-cell components are:
\begin{itemize}
\item \texttt{StochasticReactionRule} and \texttt{StochasticReactionComponent}
      for Gillespie-style SSA updates;
\item \texttt{StochasticTranscription} and \texttt{StochasticTranslation} for
      tau-leaping expression;
\item \texttt{CellGrowthComponent} for growth and density-aware mass change;
\item \texttt{CellCycleCheckpointComponent} for lifecycle phase annotation and
      clock advancement;
\item \texttt{CellFateDecisionComponent} and \texttt{CellDivisionEvent} for
      viability/division routing;
\item \texttt{PathwayStressResponseComponent} for sustained low-pathway stress;
\item \texttt{BioStateProjectionComponent} and
      \texttt{MetabolicStateProjectionComponent} for moving biochemical outputs
      into whole-cell state;
\item \texttt{CompartmentExchangeComponent} for transport between regions;
\item \texttt{EukaryoticCellCycleComponent} and \texttt{FtsZPolymerizationComponent}
      for the more specific eukaryotic and division-oriented examples;
\item \texttt{MetabolicSubmodelComponent} for compact metabolism examples;
\item \texttt{ResourceAllocationComponent} for resource partitioning examples.
\end{itemize}

The components are intentionally explicit: they share one state object, but they
do not silently merge semantics. That is why the examples can explain how a
biochemical output becomes a resource budget, a growth driver, or a lifecycle
checkpoint without pretending that these are the same thing.

\begin{figure}[htbp]
\centering
% FIGURE-SPEC-BEGIN
% ID: fig:whole-cell-architecture
% Source of truth:
%   - ../source/plugins/data/WholeCellModeling/WholeCellState.h
%   - ../source/plugins/data/WholeCellModeling/StochasticReactionRule.h
%   - ../source/plugins/components/WholeCellModeling/StochasticTranscription.h
%   - ../source/plugins/components/WholeCellModeling/StochasticTranslation.h
%   - ../source/plugins/components/WholeCellModeling/CellGrowthComponent.h
%   - ../source/plugins/components/WholeCellModeling/CellCycleCheckpointComponent.h
%   - ../source/plugins/components/WholeCellModeling/CellFateDecisionComponent.h
%   - ../source/plugins/components/WholeCellModeling/CellDivisionEvent.h
% Purpose:
%   Show the current whole-cell pipeline around the shared state object.
% Required visual content:
%   - shared state
%   - expression and degradation
%   - growth and checkpoints
%   - fate/division routing
% Accessibility description:
%   A layered pipeline where biochemical, expression, growth, and routing components all read and write the same shared state.
% Validation:
%   The diagram must match the component ordering used by the didactic whole-cell applications.
% FIGURE-SPEC-END
\figureplaceholder{figs/generated/whole-cell-architecture}
\caption{Current whole-cell architecture around \texttt{WholeCellState}.}
\label{fig:whole-cell-architecture}
\end{figure}

\section{Didactic examples}

The model-specific whole-cell examples are the clearest evidence that the
surface exists in executable form. They live under
\texttt{source/applications/modelSpecific/wcm/} and persist matching
\texttt{models/*.gen} snapshots.

The current families are:
\begin{itemize}
\item \texttt{WcmSimpleSsa}: one-gene SSA example with transcription,
      degradation, translation, and protein turnover;
\item \texttt{WcmGeneExpression}: central-dogma example with four genes and
      stochastic transcription, translation, and degradation;
\item \texttt{WcmBioNetworkBridge}: minimal deterministic BioNetwork to
      stochastic whole-cell bridge;
\item \texttt{WcmWholeCellMvp}: compact whole-cell analog combining
      deterministic metabolism, projection, expression, growth, and division;
\item \texttt{WcmFluxBalanceWholeCellMvp}: FBA-driven whole-cell MVP;
\item \texttt{WcmCompartmentTransportWholeCellMvp}: FBA plus compartment-aware
      transport and exchange;
\item \texttt{WcmPathwayStressWholeCellMvp}: stress routing when the monitored
      pathway stays below threshold;
\item \texttt{WcmLifecycleStress}: starvation and death routing example;
\item \texttt{WcmMgenitaliumKarr2012}: M. genitalium analog inspired by Karr et
      al. 2012, with explicit note that it is a didactic analog, not a canonical
      validated public whole-cell model;
\item \texttt{WcmYeastDidacticWholeCellMvp}: yeast-oriented didactic FBA/whole-
      cell example;
\item \texttt{WcmYeastEukaryoticWholeCellMvp}: budding-yeast-inspired eukaryotic
      whole-cell example with compartments and cell-cycle coordination.
\end{itemize}

The yeast examples are the most useful way to explain the architecture to a
reader. They combine FBA-style metabolism, compartment transport, expression,
growth, and lifecycle control in one didactic pipeline without presenting the
result as a validated public yeast whole-cell model.

\begin{figure}[htbp]
\centering
% FIGURE-SPEC-BEGIN
% ID: fig:yeast-mvp-structure
% Source of truth:
%   - ../source/applications/modelSpecific/wcm/WcmYeastDidacticWholeCellMvp.h
%   - ../source/applications/modelSpecific/wcm/WcmYeastEukaryoticWholeCellMvp.h
%   - ../source/applications/modelSpecific/wcm/WcmCompartmentTransportWholeCellMvp.h
%   - ../source/applications/modelSpecific/wcm/WcmPathwayStressWholeCellMvp.h
%   - ../source/applications/modelSpecific/wcm/WcmLifecycleStress.h
% Purpose:
%   Show how the yeast-oriented examples vary around one shared architectural core.
% Required visual content:
%   - common core
%   - metabolic branch
%   - compartment branch
%   - stress branch
%   - lifecycle branch
% Accessibility description:
%   A matrix or branching diagram that compares the yeast-oriented whole-cell MVP variants.
% Validation:
%   The diagram must remain aligned with the actual example class responsibilities and descriptions.
% FIGURE-SPEC-END
\figureplaceholder{figs/generated/yeast-mvp-structure}
\caption{Yeast-oriented whole-cell example structure.}
\label{fig:yeast-mvp-structure}
\end{figure}

\begin{figure}[htbp]
\centering
% FIGURE-SPEC-BEGIN
% ID: fig:pathway-stress-example
% Source of truth:
%   - ../source/applications/modelSpecific/wcm/WcmPathwayStressWholeCellMvp.h
%   - ../source/plugins/components/WholeCellModeling/PathwayStressResponseComponent.h
%   - ../source/plugins/components/WholeCellModeling/CellFateDecisionComponent.h
% Purpose:
%   Show the sustained low-pathway-stress path from objective to arrest/death routing.
% Required visual content:
%   - monitored pathway objective
%   - threshold
%   - stress streak
%   - arrest/death routing
% Accessibility description:
%   A simple lifecycle diagram where low objective activity produces arrest and death paths.
% Validation:
%   The diagram must match the current stress component inputs and routing behavior.
% FIGURE-SPEC-END
\figureplaceholder{figs/generated/pathway-stress-example}
\caption{Pathway stress example in the current whole-cell surface.}
\label{fig:pathway-stress-example}
\end{figure}

\section{Evidence and validation}

The current whole-cell and biochemical code is not only descriptive; it is also
covered by focused unit tests. The most relevant current test files are:
\begin{itemize}
\item \texttt{../source/tests/unit/test\_wcm\_plugins.cpp};
\item \texttt{../source/tests/unit/test\_simulator\_runtime.cpp}.
\end{itemize}

Those tests confirm, among other things, that:
\begin{itemize}
\item \texttt{WholeCellState} stores and returns molecule counts, metabolite
      pools, compartment pools, pathway activities, resource budgets, cell
      geometry, and lifecycle metadata;
\item \texttt{StochasticReactionRule} handles zero-rate, zero-count, unimolecular,
      and bimolecular propensity cases;
\item \texttt{BioNetwork} and \texttt{BioSimulatorRunner} expose the current
      deterministic runner and SBML import/export surface;
\item \texttt{MetabolicStateProjectionComponent} and
      \texttt{CompartmentExchangeComponent} participate in executable
      whole-cell pipelines;
\item \texttt{CellCycleCheckpointComponent} participates in executable
      whole-cell pipelines;
\item \texttt{CellFateDecisionComponent} participates in executable whole-cell
      pipelines;
\item \texttt{PathwayStressResponseComponent} participates in executable
      whole-cell pipelines;
\item \texttt{EukaryoticCellCycleComponent} participates in executable whole-
      cell pipelines;
\item SBML import/export currently round-trips the native biochemical
      representation through the implemented bridge behavior.
\end{itemize}

Figure~\ref{fig:evidence-ladder} summarizes the evidence ladder that should be
used when reading this chapter.

\begin{figure}[htbp]
\centering
% FIGURE-SPEC-BEGIN
% ID: fig:evidence-ladder
% Source of truth:
%   - ../docs/ai_assistants/GOVERNANCE.md
%   - ../docs/ai_assistants/STATUS.md
%   - ../docs/ai_assistants/reference/SCIENTIFIC_DOMAINS.md
%   - ../source/tests/unit/test_wcm_plugins.cpp
%   - ../source/tests/unit/test_simulator_runtime.cpp
% Purpose:
%   Show how code, tests, didactic examples, and scientific claims relate to each other.
% Required visual content:
%   - confirmed code
%   - executed tests
%   - didactic examples
%   - experimental/research boundary
% Accessibility description:
%   A vertical or stepped ladder that distinguishes code evidence from scientific claim level.
% Validation:
%   The diagram must not collapse experimental examples into scientific validation.
% FIGURE-SPEC-END
\figureplaceholder{figs/generated/evidence-ladder}
\caption{Evidence ladder for the biological and whole-cell surface.}
\label{fig:evidence-ladder}
\end{figure}

\section{Scientific limits and future work}

The current surface is useful, but it is not a validated organism model.
Specific limits matter:
\begin{itemize}
\item \texttt{WcmMgenitaliumKarr2012} is an analog inspired by the Karr et al.
      2012 whole-cell paper, not a canonical public model with full scientific
      validation in this repository;
\item the yeast examples are didactic architecture examples, not published
      validated yeast whole-cell simulations;
\item \texttt{HUM-VC-001} in the human backlog still owns the choice of first
      bounded virtual-cell use case;
\item \texttt{HUM-VC-002} still owns the SBML supported subset and compatibility
      target;
\item any future scientific claim must include a real reference, dataset or
      model identifier, benchmark, tolerance, procedure, and explicit maturity
      statement.
\end{itemize}

The practical takeaway is simple: use this chapter to locate the current
implementation surface and its validation boundary, not to infer a stronger
scientific claim than the code and tests support.
